%--------------------
\documentclass{article}

%-------------------------
% set input path (so that LaTeX can find the .trsl files)
\makeatletter
\def\input@path{{../}}
\makeatother
%-------------------------

% layout
\setlength{\parindent}{0pt}
\setlength{\parskip}{2ex}
%\usepackage[margin=20mm,a4paper]{geometry}
\usepackage[margin=20mm,a4paper, landscape]{geometry}
%\usepackage[margin=20mm,a3paper,landscape]{geometry}

% multinotes
\usepackage{multinotes}

% define languages
\definelanguage{welsh}{cy}{cym}{cymraeg}
\definelanguage{english}{en}{eng}{english}
\definelanguage{french}{fr}{fra}{francais}
\definelanguage{breton}{br}{bre}{brezhoneg}
\definelanguage{russian}{ru}{rus}{russkiy}
\definelanguage{thai}{th}{tha}{thai}

% include languages
\includelanguage{welsh}
\includelanguage{english}
\includelanguage{breton}
\includelanguage{french}
%\includelanguage{russian}
%\includelanguage{thai}

% enable non-latin fonts for Russian and Thai
% Note: these need LuaLaTeX (not PDFLatex)
%\babelfont{rm}{FreeSerif}
%\babelfont{sf}{FreeSans}
%\babelfont{tt}{FreeMono}

% language styles
\definecolor{darkblue}{RGB}{0,0,139}
\definecolor{darkgreen}{RGB}{2,90,2}
\definecolor{darkred}{RGB}{139,0,0}
\setlanguagestyle{welsh}{\color{red}}
\setlanguagestyle{english}{\color{blue}}
\setlanguagestyle{breton}{\color{purple}}
\setlanguagestyle{french}{\color{darkblue}}
\setlanguagestyle{russian}{\color{darkred}}
\setlanguagestyle{thai}{\color{darkgreen}}


% multicol options
\columnsep=0.05\textwidth
\columnseprule=0.4pt
\raggedright
\localcounter{section}
\localcounter{enumi}

% theorems
\usepackage{amsthm}
\theoremstyle{definition}
\newtheorem{theorem}{\GetTranslation{Theorem}}
\newtheorem{definition}[theorem]{\GetTranslation{Definition}}
\newtheorem{corollary}[theorem]{\GetTranslation{Corollary}}

\usepackage{lipsum}

%--------------------
\begin{document}

%--------------------
% title
\begin{multicols}
\cym{\title{Dogfen brawf}\author{DE}\date{\today}\maketitle}
\eng{\title{Test document}\author{DE}\date{\today}\maketitle}
\fra{\title{Document de test}\author{DE}\date{\today}\maketitle}
\bre{\title{Teul amprouiñ}\author{DE}\date{\today}\maketitle}
\rus{\title{Тестовый документ}\author{DE}\date{\today}\maketitle}
\tha{\title{เอกสารทดสอบ}\author{DE}\date{\today}\maketitle}
\end{multicols}


% abstract
\begin{multicols}

\begin{cymraeg}
\begin{abstract}
Helo byd!
\end{abstract}
\end{cymraeg}

\begin{english}
\begin{abstract}
Hello world!
\end{abstract}
\end{english}

\begin{brezhoneg}
\begin{abstract}
Demat bed!
\end{abstract}
\end{brezhoneg}

\begin{francais}
\begin{abstract}
Bonjour le monde!
\end{abstract}
\end{francais}

\begin{russkiy}
\begin{abstract}
Привет, мир!
\end{abstract}
\end{russkiy}

\begin{thai}
\begin{abstract}
สวัสดีโลก
\end{abstract}
\end{thai}

\end{multicols}

%--------------------
\section{Introduction}
Some common text
\cym{\lipsum[1]}
\eng{\lipsum[2]}
\bre{\lipsum[3]}
\fra{\lipsum[4]}
\rus{\lipsum[5]}
\tha{\lipsum[6]}

%--------------------
\section{Sequential mode}

\begin{cymraeg}
\begin{definition}
\lipsum[1]
\end{definition}
\end{cymraeg}

\begin{english}
\begin{definition}
\lipsum[2]
\end{definition}
\end{english}

\begin{brezhoneg}
\begin{definition}
\lipsum[3]
\end{definition}
\end{brezhoneg}

\begin{francais}
\begin{definition}
\lipsum[4]
\end{definition}
\end{francais}

\begin{russkiy}
\begin{definition}
Это простой, фиктивный контент, используемый в типографском или наборном бизнесе.
\end{definition}
\end{russkiy}

\begin{thai}
\begin{definition}
คือ เนื้อหาจำลองแบบเรียบๆ ที่ใช้กันในธุรกิจงานพิมพ์หรืองานเรียงพิมพ์
\end{definition}
\end{thai}

%--------------------
\section{Parallel mode}

\begin{multicols}

\begin{english}
\begin{definition}
\lipsum[1]
\end{definition}
\end{english}

\begin{cymraeg}
\begin{definition}
\lipsum[2]
\end{definition}
\end{cymraeg}

\begin{brezhoneg}
\begin{definition}
\lipsum[3]
\end{definition}
\end{brezhoneg}

\begin{francais}
\begin{definition}
\lipsum[4]
\end{definition}
\end{francais}

\begin{russkiy}
\begin{definition}
\lipsum[5]
\end{definition}
\end{russkiy}

\begin{thai}
\begin{definition}
\lipsum[6]
\end{definition}
\end{thai}

\end{multicols}

%--------------------
\section{Parallel mode using 
{\tt\textbackslash multicolsync}
and
{\tt\textbackslash multicolspan}}

%\multicolsoff

\begin{multicols}

\cym{\lipsum[1]}
\eng{\lipsum[2]}
\bre{\lipsum[3]}
\fra{\lipsum[4]}
\rus{\lipsum[5]}
\tha{\lipsum[6]}

\multicolsync

\cym{\lipsum[7]}
\eng{\lipsum[8]}
\bre{\lipsum[9]}
\fra{\lipsum[10]}
\rus{\lipsum[11]}
\tha{\lipsum[12]}

\multicolspan{\lipsum[13]}

\cym{\lipsum[14]}
\eng{\lipsum[15]}
\bre{\lipsum[16]}
\fra{\lipsum[17]}
\rus{\lipsum[18]}
\tha{\lipsum[19]}

\end{multicols}

\end{document}